\section{The remaining route to general quasi-polynomial time}
\label{sec:route}

\subsection{Two routes with different missing theorems}
The current implementation can now avoid work in three distinct ways:
recognize a certified geometric class immediately; share repeated direct
summands; or prove sufficient final homology dimension from a partial scan.
None of these mechanisms gives a universal estimate on the frontier or on
the connected differential components.

The algebraic route needs a theorem that constructs a useful decomposition
of \emph{every} knot diagram, with controlled component complexity under
continuation. The hierarchy route needs an executable representation of
patterned handle structures, compressed surface operations, and a global
bound on effective restarts. These are different programmes. A successful
subroutine for one programme does not automatically discharge the other
programme's complexity obligations.

For the algebraic route, a quantitative target is to prove for some
constructively specified ordering or more general assembly rule that
\[
 b^2 2^{w/2}=O(\log^2 n).
\]
This target is sufficient rather than necessary. The type-count proof is
crude enough that a much sharper classification of reachable morphisms or
templates could replace it. It may be more realistic to control the number
of \emph{reachable} templates directly than to force every template to
have logarithmic size.

There is also a topological obstruction to reducing the frontier on every
input merely by choosing a better crossing order. The tree-width theorem
of de Mesmay, Purcell, Schleimer, and Sedgwick implies that every diagram
of a torus knot $T(p,q)$ has tree-width at least
$\min(p,q)/4-1$~\cite[Theorem 3.7]{treewidth}. A linear frontier scan
induces a carving decomposition up to constant factors, so the standard
$n=\Theta(p^2)$ diagrams of $T(p,p+1)$ require $w=\Omega(\sqrt n)$
under every such ordering. Their homogeneous structure is already handled
by the new front end. A useful general width theorem must therefore concern
the residual class after certificates, or a genuinely different
representation. Width of a knot diagram and width of a triangulated knot
complement must not be conflated; the same source shows that torus-knot
complements admit constant-tree-width triangulations.

\subsection{Nonuniform restart accounting for a future hierarchy engine}
The following elementary result is useful when different hierarchy levels
have substantially different complexity caps. It is stated independently of
any claim that the necessary geometric invariants have been implemented.

\begin{proposition}[Mixed-radix restart bound]\label{prop:radix}
Suppose an algorithm has a length-$L$ state vector
$(a_1,\ldots,a_L)$ with integer bounds $0\le a_i\le g_i$. At each restart,
the first changed coordinate strictly decreases; all subsequent coordinates
may reset arbitrarily within their fixed bounds. Then the number of such
restarts is at most
\[
 \prod_{i=1}^{L}(g_i+1)-1.
\]
If every between-restart operation has polynomial bit cost and all stored
encodings have polynomial size, a sufficient condition for the target
$n^{O(\log n)}$ is
\[
 \sum_{i=1}^{L}\log(g_i+1)=O(\log^2 n),
\]
with polynomially many operations between successive restarts.
\end{proposition}
\begin{proof}
Use the mixed-radix potential
\[
 \Phi(a_1,\ldots,a_L)=\sum_{i=1}^{L}a_i\prod_{j>i}(g_j+1).
\]
It is an integer between zero and $\prod_i(g_i+1)-1$. If coordinate $i$
is the first to change, decreasing it by one lowers $\Phi$ by
$\prod_{j>i}(g_j+1)$. The entire suffix can increase by at most that
quantity minus one. Every restart therefore decreases $\Phi$ by at least
one. Taking logarithms of the resulting product bound gives the last
statement after charging the polynomial overhead.
\end{proof}

The uniform case with $g_i=g$ gives the familiar $(g+1)^L$ factor. Polynomial
surface caps with $L=O(n)$ still permit $n^{O(n)}$ restarts. Polynomial caps
with $L=O(\log n)$ yield~\eqref{eq:target}; with
$L=O((\log n)^2)$ they give $2^{O((\log n)^3)}$. Coordinate compression
alone does not change this arithmetic.

The proposition is a template for a proof, not an observation that the
geometric algorithm already has this potential. The caps must be globally
valid after every reconstruction. A log recording decreasing numbers does
not establish either that validity or the topological correctness of the
restart operation.

\subsection{Concrete interface obligations for a hierarchy implementation}
The minimum coherent implementation unit is larger than a normal-coordinate
vector or a ball-pattern tester. It should include the following contracts.
\begin{enumerate}
\item \textbf{Encoded topology.} A finite list of allowed handle types, exact
attachment maps, boundary-pattern incidences, orientation, and an encoding
size charged to the input and all created objects.
\item \textbf{Compressed surface operations.} Binary multiplicities,
component and boundary information, cutting, gluing, and transport of curves
or discs across each change. A procedure must operate on binary size, rather
than expand one disc per unit of normal weight.
\item \textbf{Topology-preserving simplification.} Every parallelity-bundle
replacement or handle cleanup must carry a precise statement of the
manifold and boundary pattern it preserves or changes.
\item \textbf{Progress.} Surface production and violating-disc transport
must support a globally bounded potential through every replacement of a
hierarchy suffix, including changes to the representation itself.
\item \textbf{Cost composition.} Bound the number of generated handles,
coefficient lengths, calls, and restarts together. A fast primitive is
insufficient if called super-quasi-polynomially often.
\end{enumerate}
The Agol--Hass--Thurston orbit-count method is an established model for
processing exponentially long interval systems through short encodings
\cite{aht}. Its existence supports this design direction; it does not, by
itself, implement the full contracts above. The present archive leaves
these interfaces as a research programme rather than supplying placeholders
that report topological verdicts.

\section{Further research questions and proposed work}
\label{sec:questions}

The following questions are ordered by the strength and immediacy of the
connection to the delivered implementation. Each is paired with a concrete
success criterion.

\begin{enumerate}[label=\textbf{Q\arabic*.},leftmargin=*]
\item \textbf{Which prime families have bounded component size after sharing?}
Find an infinite family of prime, nonhomogeneous diagrams for which the
actual component-sharing policy has a proved bound on $b$ and $w$, and for
which earlier filters are not the entire explanation of fast recognition.
A proof must specify the scan order and the representative transition types.
A list of observed small components is useful evidence but not the result.

\item \textbf{Can reachable morphisms replace the all-masks bound?}
The factor $2^{s^2 2^{w/2}}$ counts every matrix entry in the full dot-mask
space. Identify a closure property of the entries actually produced by
extension and unit elimination. A bound on the number of reachable
coefficient masks could enlarge the quasi-polynomial regime substantially.
The proposed class must survive composition and Schur-complement addition.

\item \textbf{Can connected components be split efficiently after a basis
change?} The current graph decomposition finds literal blocks. Develop a
restricted, provably polynomial family of basis changes that reveals further
repeated summands. The output should include an explicit invertible change
of basis or chain-homotopy equivalence. A general canonical-form oracle
would merely move the hard computation elsewhere.

\item \textbf{How much can a pointed or reduced category expose?}
A carefully constructed marked-edge implementation might kill some
nonunit morphisms earlier and reveal smaller components. It must track the
basepoint through all gluing and delooping operations, prove compatibility
with reduced homology, and preserve a finite composition algebra. This
archive does not include an unverified quotient obtained by simply deleting
dot terms touching a chosen boundary point.

\item \textbf{Can continuation-aware obstructions be strengthened?}
Theorem~\ref{thm:euler} uses one scalar Euler functional for each completed
summand. Investigate several quantum-graded evaluations or other additive
functionals with certified norm bounds. The central requirement is a proved
lower bound on homology rank after continuation. Capping signed sums before
cancellation would invalidate such a certificate.

\item \textbf{When should the optional Euler engine run?}
Develop a predictor based on cheap component statistics or a small number
of exact suffix values. Measure total time including failed inference.
Any adaptive policy must preserve the property that a budget failure simply
returns to a complete backend. Establishing a profitable policy on prime
filter-resistant diagrams would be more informative than optimizing the
visible-sum stress family again.

\item \textbf{Can a complete three-braid rank engine be integrated?}
The polynomial-time three-braid result in the 2026 literature suggests a
specialized exact backend based on normal forms and formulas
\cite{kelomaki}. Implementation requires careful grading, mirror, linking,
and coefficient conventions. Validate it against the independent cube and
both scanning algebras before using it to bypass scanning. The mere fact
that the input has three strands is not a complexity theorem for the
current scanner.

\item \textbf{Can extremal-band complexity be made uniform?}
Seek a fully quantified finite-field bound such as
$2^{O(w)}n^{O(k)}$ for computing a band of $k$ homological degrees under an
explicitly specified assembly rule. The number of objects, decorated
morphisms, coefficient sizes, and the extra degree needed for outgoing
differentials must all be charged. A complete recognizer would still need a
proved location of a nontriviality witness in those bands or a fallback.

\item \textbf{What structural parameter survives reduction?}
Example~\ref{ex:defectone} and the I/II-reduced family in
Section~\ref{sec:reduced-defect-one} rule out control by homogeneity defect
alone. Search for a parameter combining signed block structure with a
certified diagram reduction or tangle decomposition. It should distinguish
many redundant twists from genuine mixed-sign complexity.

\item \textbf{Can a controlled portfolio improve the heavy tail?}
Sharing, saturated sharing, Euler inference, the ordinary scanner, and
geometric reduction have different overheads. A portfolio should be judged
by CPU work as well as wall time. A useful guarantee would compare its cost
with the best applicable complete backend, while bounding the work lost in
abandoned runs. Purely retrospective selection of the fastest run is not an
implementable performance claim.

\item \textbf{Can the new proof obligations be formalized?}
A small formal development can first prove the $S_3$ semiring quotient,
weighted direct-sum transition invariant, and monotone Euler bound. These
have explicit finite algebraic interfaces. Formalizing the full
Bar-Natan category and the unknot-detection theorem is a separate,
substantially larger task. A checked combinatorial certificate still depends
on its imported topological soundness theorem.

\item \textbf{Can the hierarchy route obtain a nonuniform progress bound?}
Implement a complete unaccelerated hierarchy engine with explicit contracts,
then search for level-dependent caps satisfying
Proposition~\ref{prop:radix}. A successful result would specify the caps,
the retained suffix under each simplification, and the proof that every
expensive restart is counted. This is a direct path to a real end-to-end
quasi-polynomial theorem, but it remains beyond the delivered changes.
\end{enumerate}

\section{Concluding assessment}
The delivered recognizer has a broader polynomially decidable front end and
an exact mechanism for representing repeated chain-complex summands without
expanding their generators. Saturated weights specialize that representation
to the actual yes/no problem, and exact Euler continuations supply a safe
form of early nontriviality detection. The finite-type theorem identifies a
specific sufficient regime for $n^{O(\log n)}$ time that accounts for both
algebraic state size and integer arithmetic.

The implementation retains an exponential worst-case upper bound; no general
quasi-polynomial bound is proved. The
measured improvement on homogeneous families is an end-to-end pipeline
improvement after validation. The dramatic connected-sum scanner figures
are representation stress tests, because an earlier specialized factorizer
already handles those examples efficiently. The hardest remaining task is
to control structural complexity for diagrams that are neither certified by
the front end nor decomposed into small repeated components. The archive
provides an executable basis for pursuing that task and exact criteria for
judging whether a future claim has actually reached the quasi-polynomial
target.
